\chapter{Volatilité implicite et sa surface}\label{ch:dv:implied-volatility-and-its-surface}

Avant octobre 1987, les volatilités implicites des options sur l'indice S\&P~500 étaient
à peu près identiques à chaque strike~: une droite plate, telle que la prévoit le \omterm{def:dv:black-scholes-three-ways:model}{modèle
de Black--Scholes}. Après le krach, au cours duquel l'indice a perdu un cinquième en deux
jours, les strikes bas ont commencé à se traiter à des volatilités plus élevées que les
strikes hauts, et ils le font depuis lors. Aujourd'hui, un jour calme, un put à trois mois
10~\% sous le marché se traite à une volatilité de 20, tandis que le call situé 10~\% au-dessus se traite à 14~; un put
20~\% sous le marché se traite à 27, et valorise une baisse que la volatilité à la monnaie de 16
qualifierait d'événement d'une chance sur trois cents, et que le prix traite comme un
événement d'une chance sur cinquante. Le modèle n'a pas changé~; ce sont les prix qui ont changé. Depuis
lors, le modèle sert d'unité de mesure~: chaque option est cotée par la volatilité que la formule
requiert pour reproduire son prix, et le tableau de ces nombres selon les strikes et les
échéances est ce que regarde, ajuste, négocie et gère en risque un desk d'options. Ce chapitre
décrit ce tableau, les règles qu'il doit respecter, et ce qu'il dit de la
distribution que le marché valorise.

\section{La surface comme liste de prix du marché}

La volatilité implicite est définie dans le livre~1, chapitre~25, avec le skew
selon les strikes et la structure par terme selon les échéances. Réunis, ils forment une
surface, et deux changements de coordonnées la rendent comparable d'une échéance à l'autre et
d'un sous-jacent à l'autre.

\begin{definition}[Surface de volatilité, log-moneyness, variance implicite totale]\label{def:dv:implied-volatility-and-its-surface:surface}
La \emph{surface de volatilité}\index{surface de volatilité} d'un sous-jacent est la
fonction $(K,T)\mapsto\sigma_{\mathrm{imp}}(K,T)$ qui donne, pour chaque strike et
échéance, la volatilité implicite de l'option européenne, calculée avec la formule de Black
sur le forward $F_{0,T}$ de cette échéance. Ses coordonnées naturelles sont la
\emph{log-moneyness}\index{log-moneyness} $k=\ln(K/F_{0,T})$, qui mesure un strike
par rapport au forward de sa propre échéance, et la \emph{variance implicite
totale}\index{variance implicite totale} $w(k,T)=\sigma_{\mathrm{imp}}^2T$, la variance
du log-rendement à l'échéance que le prix implique.
\end{definition}

Dans ces coordonnées, la formule de Black ne dépend que de $(k,w)$, à un facteur près
$P(0,T)F_{0,T}$~: $C/(PF)=\Phi(-k/\sqrt w+\sqrt w/2)-e^k\Phi(-k/\sqrt w-\sqrt w/2)$.
Une surface est alors une famille de smiles $k\mapsto w(k,T)$, une par échéance, et les
conditions d'absence d'\omterm{def:dv:implied-volatility-and-its-surface:static}{arbitrage statique} des sections suivantes deviennent des conditions sur $w$.

\begin{omfigure}
\begin{tikzpicture}
\begin{axis}[width=10.5cm,height=5cm,
  xlabel={prix d'exercice (spot 100)},ylabel={volatilité implicite (\%)},
  tick label style={font=\small},label style={font=\small},
  xmin=60,xmax=130,ymin=8,ymax=55,grid=major,grid style={gray!25},
  legend columns=4,legend style={font=\footnotesize,at={(0.5,-0.3)},anchor=north,draw=none,fill=none,/tikz/every even column/.append style={column sep=8pt}},
  table/col sep=comma]
\addplot[omProp,thick] table[x=strike,y=m1]{figdata/derivatives/07-implied-volatility-and-its-surface/smiles.csv};
\addplot[omDef,thick] table[x=strike,y=m3]{figdata/derivatives/07-implied-volatility-and-its-surface/smiles.csv};
\addplot[omThm,very thick] table[x=strike,y=y1]{figdata/derivatives/07-implied-volatility-and-its-surface/smiles.csv};
\addplot[gray,very thick,dashed] table[x=strike,y=y2]{figdata/derivatives/07-implied-volatility-and-its-surface/smiles.csv};
\legend{1 mois,3 mois,1 an,2 ans}
\end{axis}
\end{tikzpicture}

\omcaption{La \omterm{def:dv:implied-volatility-and-its-surface:surface}{surface de volatilité} des indices actions du chapitre~: quatre smiles d'une même surface. Les échéances courtes
sont pentues et basses à la monnaie~; les échéances longues sont plus plates et plus hautes. Le
skew 90--110 vaut 10.8 points de volatilité à un mois, 6.3 à trois mois, 3.3 à un an
et 2.5 à deux ans. Paramètres illustratifs, en accord avec les formes décrites dans
le texte.}\label{fig:dv:implied-volatility-and-its-surface:smiles}
\end{omfigure}

La surface est aussi la liste de prix du desk. Une cotation de toute option européenne sur le
sous-jacent, aussi inhabituels soient son strike ou son échéance, se lit sur elle~; un produit structuré
évalué par un modèle calibré sur elle hérite de son skew~; et le risque de l'ensemble du
portefeuille se mesure face à ses mouvements.

\section{Phénoménologie des smiles par classe d'actifs}

Les formes diffèrent d'un marché à l'autre, et ces différences ont des raisons.

\begin{itemize}
\item \textbf{Indices actions}~: un skew négatif prononcé, plus marqué pour les échéances courtes et
décroissant à peu près comme une puissance de l'échéance (chapitre~12), avec un niveau
à la monnaie bas sur les marchés calmes et une structure par terme croissante. Les investisseurs achètent des puts sur indices comme
une assurance, les krachs se produisent à la baisse, et la volatilité monte quand les cours baissent.
\item \textbf{Actions individuelles}~: un skew de même sens mais plus faible, avec davantage de courbure~: les études
portant sur les options d'actions individuelles montrent que leurs distributions risque-neutres sont
nettement moins asymétriques vers la baisse que celle de l'indice. Le skew d'indice ajoute la tendance des actions à baisser ensemble,
soit le skew de corrélation du chapitre~17.
\item \textbf{Devises}~: des smiles presque symétriques autour du niveau à la monnaie, inclinés
par le risk reversal vers la devise que l'on craint de voir baisser (livre~2,
chapitre~19).
\item \textbf{Matières premières}~: souvent asymétriques vers le haut. Un choc d'offre fait monter le prix et la
volatilité ensemble, et la plupart des matières premières affichent une volatilité plus élevée après
les hausses de prix, à l'inverse du schéma des actions~; le pétrole brut constitue l'exception
documentée (\cref{fig:dv:implied-volatility-and-its-surface:classes}).
\item \textbf{Taux d'intérêt}~: cotés en volatilité normale (livre~2, chapitre~13),
avec des smiles modélisés par SABR (chapitre~11) et un cube d'échéances, de maturités et de prix d'exercice
construit dans le livre~6, chapitre~5.
\end{itemize}

\begin{omfigure}
\begin{tikzpicture}
\begin{axis}[width=10.5cm,height=5cm,
  xlabel={log-moneyness $k$},ylabel={volatilité implicite (\%)},
  tick label style={font=\small},label style={font=\small},
  xmin=-0.3,xmax=0.3,ymin=5,ymax=42,grid=major,grid style={gray!25},
  legend columns=3,legend style={font=\footnotesize,at={(0.5,-0.3)},anchor=north,draw=none,fill=none,/tikz/every even column/.append style={column sep=8pt}},
  table/col sep=comma]
\addplot[omThm,very thick] table[x=k,y=equity]{figdata/derivatives/07-implied-volatility-and-its-surface/asset_classes.csv};
\addplot[omDef,very thick] table[x=k,y=fx]{figdata/derivatives/07-implied-volatility-and-its-surface/asset_classes.csv};
\addplot[omProp,very thick,dashed] table[x=k,y=commodity]{figdata/derivatives/07-implied-volatility-and-its-surface/asset_classes.csv};
\legend{indice actions,paire de devises,matière première}
\end{axis}
\end{tikzpicture}

\omcaption{Smiles à trois mois par classe d'actifs, en fonction de la \omterm{def:dv:implied-volatility-and-its-surface:surface}{log-moneyness}~: un indice
actions asymétrique vers le bas, une paire de devises presque symétrique, une matière première asymétrique vers le haut.
Formes illustratives.}\label{fig:dv:implied-volatility-and-its-surface:classes}
\end{omfigure}

\section{Sticky strike, sticky delta et ce que suppose le desk}

Une surface est un instantané. Lorsque le spot se déplace, il faut faire une hypothèse sur la
manière dont la surface se déplace avec lui, et cette hypothèse modifie la couverture.

\begin{definition}[Sticky strike et sticky delta]\label{def:dv:implied-volatility-and-its-surface:sticky}
Dans le régime \emph{sticky strike}\index{sticky strike}, la volatilité implicite de chaque strike
et de chaque échéance reste fixe lorsque le spot se déplace~: le smile, tracé en fonction du strike,
ne bouge pas. Dans le régime \emph{sticky delta}\index{sticky delta} (aussi appelé sticky moneyness),
la volatilité implicite de chaque moneyness $K/F$ reste fixe~: le smile suit le forward, et
la volatilité à la monnaie ne change pas.
\end{definition}

Les deux règles sont des régimes plutôt que des lois~: les études sur les volatilités des indices trouvent des périodes
où chacune décrit mieux le marché, et une troisième dans laquelle la volatilité à la monnaie varie davantage
que ne le prédit l'une ou l'autre (Derman). Elles importent pour trois raisons.
Elles donnent des P\&L différents pour un même mouvement~: après une baisse de 5~\% de l'indice, le
put à trois mois de strike 100 vaut 6.00 sous \omterm{def:dv:implied-volatility-and-its-surface:sticky}{sticky strike} et 5.78 sous \omterm{def:dv:implied-volatility-and-its-surface:sticky}{sticky
delta}, contre 3.08 avant (\cref{fig:dv:implied-volatility-and-its-surface:regimes}).
Elles donnent aussi des deltas différents.

\begin{proposition}[Delta ajusté au smile]\label{prop:dv:implied-volatility-and-its-surface:delta}
Si la volatilité implicite d'un strike fixe varie avec le spot selon $\partial\sigma_{\mathrm{imp}}/\partial S$,
le delta total de l'option vaut $\Delta_{\mathrm{BS}}+\mathcal V\,\partial\sigma_{\mathrm{imp}}/\partial S$.
Sous \omterm{def:dv:implied-volatility-and-its-surface:sticky}{sticky strike}, la correction est nulle~; sous \omterm{def:dv:implied-volatility-and-its-surface:sticky}{sticky delta}, le smile étant une fonction
de $k=\ln(K/F)$, $\partial\sigma_{\mathrm{imp}}/\partial S=-\sigma_{\mathrm{imp}}'(k)/S$, qui est
positif pour un skew négatif.
\end{proposition}
\begin{proof}
Dérivons $V=C_{\mathrm{BS}}(S,\sigma_{\mathrm{imp}}(K,T;S))$ par rapport à $S$~; avec
$\sigma_{\mathrm{imp}}=\Sigma(\ln K-\ln F)$ et $F$ proportionnel à $S$,
$\partial\sigma_{\mathrm{imp}}/\partial S=-\Sigma'(k)/S$.
\end{proof}

Pour le call à trois mois de strike 100 sur la surface du chapitre, le delta de Black--Scholes
vaut 0.533 et le delta \omterm{def:dv:implied-volatility-and-its-surface:sticky}{sticky delta} 0.593~: un desk qui suppose un régime et travaille dans l'autre
est sous-couvert ou sur-couvert de six actions par centaine d'options. Elles donnent aussi des
prix exotiques différents, car un modèle de la dynamique de la surface (volatilité locale, chapitre 9~;
volatilité stochastique, chapitre 10) implique l'un ou l'autre régime~; le chapitre 11 mesure
l'écart.

\begin{omfigure}
\begin{tikzpicture}
\begin{axis}[width=10.5cm,height=5cm,
  xlabel={strike},ylabel={volatilité implicite à trois mois (\%)},
  tick label style={font=\small},label style={font=\small},
  xmin=80,xmax=120,ymin=10,ymax=28,grid=major,grid style={gray!25},
  legend columns=2,legend style={font=\footnotesize,at={(0.5,-0.3)},anchor=north,draw=none,fill=none,/tikz/every even column/.append style={column sep=8pt}},
  table/col sep=comma]
\addplot[omDef,very thick] table[x=strike,y=before]{figdata/derivatives/07-implied-volatility-and-its-surface/regimes.csv};
\addplot[omThm,very thick,dashed] table[x=strike,y=sticky_delta]{figdata/derivatives/07-implied-volatility-and-its-surface/regimes.csv};
\draw[omDef] (axis cs:100.37,16.36) circle[radius=2.5pt];
\draw[omThm] (axis cs:95.36,16.36) circle[radius=2.5pt];
\draw[omDef,dotted,thick] (axis cs:95.36,10) -- (axis cs:95.36,18.08);
\node[font=\footnotesize,anchor=west,omDef,align=left] at (axis cs:96,21.5) {nouveau strike à la monnaie 95.4~:\\18.1 sous sticky strike};
\legend{{avant, et après la baisse sous sticky strike},après une baisse de 5~\% sous sticky delta}
\end{axis}
\end{tikzpicture}

\omcaption{Le smile à trois mois avant et après une baisse de 5~\% de l'indice. Sous \omterm{def:dv:implied-volatility-and-its-surface:sticky}{sticky
strike}, la courbe reste en place et le nouveau strike à la monnaie (95.4) lit une volatilité de 18.1
sur celle-ci~; sous \omterm{def:dv:implied-volatility-and-its-surface:sticky}{sticky delta}, la courbe se décale vers la gauche avec le forward et la
volatilité à la monnaie reste à 16.4. Marqueurs ronds~: points à la monnaie. Données~: le
tutoriel.}\label{fig:dv:implied-volatility-and-its-surface:regimes}
\end{omfigure}

\section{Conditions de non-arbitrage statique}

Tous les tableaux de nombres ne constituent pas une surface. Le chapitre 1 a énuméré les inégalités sans modèle sur les
prix d'une seule échéance~; à travers les échéances, une famille supplémentaire est nécessaire.

\begin{definition}[Arbitrage statique, arbitrage butterfly et arbitrage calendaire]\label{def:dv:implied-volatility-and-its-surface:static}
Un \emph{arbitrage statique}\index{arbitrage statique} dans un ensemble de cotations d'options est un portefeuille
d'options, de forward et d'obligations, acheté une fois et conservé jusqu'aux échéances, qui
ne coûte rien et verse un montant non négatif dans chaque état, strictement positif dans certains. Un
\emph{arbitrage butterfly}\index{arbitrage butterfly} est un arbitrage au sein d'une même échéance~: des prix de call
qui ne sont pas convexes en fonction du strike, de sorte qu'un butterfly a un prix négatif. Un
\emph{arbitrage calendaire}\index{arbitrage calendaire} est un arbitrage entre échéances~: une option à échéance plus lointaine
moins chère qu'une option à échéance plus proche à même moneyness forward.
\end{definition}

L'absence d'arbitrage de call spread, de butterfly et calendaire suffit à exclure tout
\omterm{def:dv:implied-volatility-and-its-surface:static}{arbitrage statique} d'une surface de cotations sur un même sous-jacent (Carr et Madan). Exprimées en
variance totale, les conditions sont simples.

\begin{proposition}[Arbitrage statique en variance totale]\label{prop:dv:implied-volatility-and-its-surface:static}
Avec des taux et des dividendes déterministes, proportionnels à l'action~: (i) il n'y a pas
d'\omterm{def:dv:implied-volatility-and-its-surface:static}{arbitrage calendaire} si et seulement si $w(k,T)$ est non décroissante en $T$ pour tout $k$~; (ii) un smile
$k\mapsto w(k)$ est exempt d'\omterm{def:dv:implied-volatility-and-its-surface:static}{arbitrage butterfly} si et seulement si les prix de call qu'il implique
sont convexes en fonction du strike, ce qui est le cas lorsque
\[
g(k)=\Bigl(1-\frac{kw'}{2w}\Bigr)^2-\frac{w'^2}{4}\Bigl(\frac1w+\frac14\Bigr)+\frac{w^{\prime\prime}}{2}\ \ge0
\]
pour tout $k$, avec des prix de call qui tendent vers zéro pour les grands strikes.
\end{proposition}
\begin{proof}[Preuve partielle]
(i) À $k$ fixé, le call normalisé $C/(PF)$ est une fonction croissante de $w$ seule~; un calendar spread
à moneyness forward fixée est un call sur une martingale à deux dates, dont la valeur croît
avec la date par l'inégalité de Jensen, donc $w$ doit croître. (ii) La dérivée seconde du call
par rapport au strike égale la densité de la section suivante, et son calcul à travers $w(k)$
donne $g(k)$ multiplié par un facteur positif $\varphi(d_2)/\sqrt w$~; l'énoncé complet
est démontré par Gatheral et Jacquier.
\end{proof}

Une seule cotation erronée suffit à briser une surface. Sur la surface du chapitre, relever
de trois points la volatilité à trois mois en $k=-0.1$ rend $g$ négatif en six points de ce smile,
jusqu'à $-2.06$~: le butterfly centré en ce point aurait un prix négatif. Abaisser de sept points
la cotation à la monnaie à six mois fait décroître $w$ entre trois et six mois~: il y a un \omterm{def:dv:implied-volatility-and-its-surface:static}{arbitrage
calendaire}. Le constructeur de surfaces de ce chapitre vérifie ces deux points avant que tout modèle ne soit
calibré sur la surface (chapitre 8).

\section{La densité risque-neutre}

\begin{definition}[Densité risque-neutre et formule de Breeden--Litzenberger]\label{def:dv:implied-volatility-and-its-surface:density}
La \emph{densité risque-neutre}\index{densité risque-neutre} de $S_T$ est la densité
$q_T(s)$ du sous-jacent en $T$ sous la probabilité risque-neutre (livre~4,
chapitre~5)~; c'est la version continue des \omterm{def:dv:no-arbitrage-and-the-fundamental-theorems:ad}{prix d'état} du chapitre~1, divisés par le
facteur d'actualisation. La \emph{formule de Breeden--Litzenberger}\index{formule de Breeden--Litzenberger}
la lit directement dans les prix des calls~:
\[
q_T(K)=\frac{1}{P(0,T)}\,\frac{\partial^2C(K,T)}{\partial K^2}.
\]
\end{definition}
\begin{proof}[Démonstration]
Le call est l'espérance actualisée du gain,
\[
C(K,T)=P(0,T)\int_K^\infty(s-K)\,q_T(s)\,ds .
\]
Une dérivation en $K$ donne $-P(0,T)\int_K^\infty q_T(s)\,ds$, soit le digital actualisé
changé de signe, et une seconde dérivation donne la formule.
\end{proof}

La dérivée première dit déjà quelque chose du skew~: le prix d'un call digital est
$-\partial C/\partial K=P(0,T)\Phi(d_2)-\mathcal V\,\partial\sigma_{\mathrm{imp}}/\partial K$,
soit le digital de Black--Scholes augmenté d'une correction de skew. Avec un skew négatif, la correction est
positive~: le call digital à la monnaie à trois mois sur la surface du chapitre vaut 0.558,
contre 0.498 sans skew.

\begin{example}[La prime de krach]\label{ex:dv:implied-volatility-and-its-surface:crash}
Sur la surface du chapitre, le forward à trois mois vaut 100.37 et la volatilité à la monnaie
16.36~\%. Avec une surface plate à ce niveau, la probabilité risque-neutre que l'indice termine
sous 80 vaut 0.31~\%~; avec le skew, dont le strike 80 se situe à 26.8 de volatilité, elle vaut 1.89~\%, soit
six fois plus (\cref{fig:dv:implied-volatility-and-its-surface:density}). Le put 80 coûte 0.006
sur la surface plate et 0.219 sur la surface avec skew.
\end{example}

\begin{omfigure}
\begin{tikzpicture}
\begin{axis}[width=10.5cm,height=5cm,
  xlabel={niveau de l'indice à 3 mois},ylabel={densité risque-neutre},
  tick label style={font=\small},label style={font=\small},
  xmin=55,xmax=140,ymin=0,ymax=0.058,grid=major,grid style={gray!25},
  yticklabel style={/pgf/number format/.cd,fixed,precision=2},scaled y ticks=false,
  legend columns=2,legend style={font=\footnotesize,at={(0.5,-0.3)},anchor=north,draw=none,fill=none,/tikz/every even column/.append style={column sep=8pt}},
  table/col sep=comma]
\addplot[omThm,very thick] table[x=s,y=skewed]{figdata/derivatives/07-implied-volatility-and-its-surface/density.csv};
\addplot[omDef,thick,dashed] table[x=s,y=flat]{figdata/derivatives/07-implied-volatility-and-its-surface/density.csv};
\draw[gray,dotted,thick] (axis cs:80,0) -- (axis cs:80,0.058);
\legend{surface avec skew,plat à la volatilité à la monnaie}
\end{axis}
\end{tikzpicture}

\omcaption{\omterm{def:dv:implied-volatility-and-its-surface:density}{Densités risque-neutres} de l'indice à trois mois, obtenues par Breeden--Litzenberger
à partir des prix des calls~: avec le skew du chapitre (trait plein) et avec une surface plate
à la même volatilité à la monnaie (tirets). Le skew déplace la probabilité vers la queue gauche et fait
que le sommet de la densité se situe au-dessus du forward~; la probabilité d'être sous 80 (pointillés) passe
de 0.31~\% à 1.89~\%. Données~: le tutoriel.}\label{fig:dv:implied-volatility-and-its-surface:density}
\end{omfigure}

La densité est un prix, non une prévision. Elle intègre ce que les investisseurs paient pour se protéger, et
sa queue gauche est plus épaisse que la fréquence réelle des krachs si l'assurance se paie une
prime, comme le mesure le chapitre~14 sur la variance. Ce qu'elle représente exactement est le portefeuille d'options
qui verse une unité si l'indice termine dans un petit intervalle autour de $s$~: le \omterm{def:dv:no-arbitrage-and-the-fundamental-theorems:ad}{titre d'Arrow--Debreu}
du chapitre~1, rendu continu et réplicable à partir du strip d'options cotées.

\section{Tutoriel~: construction et vérification d'une surface}\label{tut:dv:implied-volatility-and-its-surface:tut}

\begin{tutorial}
\textbf{Objectif.} Construire une surface à partir de cotations en variance totale, la vérifier contre l'\omterm{def:dv:implied-volatility-and-its-surface:static}{arbitrage
statique}, la casser avec une cotation erronée, et lire la \omterm{def:dv:implied-volatility-and-its-surface:density}{densité risque-neutre} qui en découle. \textbf{État final~:}
\cref{fig:dv:implied-volatility-and-its-surface:smiles,fig:dv:implied-volatility-and-its-surface:density}
et les chiffres de l'\cref{ex:dv:implied-volatility-and-its-surface:crash}.
\begin{tutsteps}
\item \textbf{Les contrôles}~: le facteur de densité $g$ sur chaque tranche, et $w$ d'une tranche à l'autre.
\omcode{code/firm/volsurface/firm_volsurface.py}{102}{120}{Contrôles de butterfly et de calendrier en variance totale.}\label{lst:dv:implied-volatility-and-its-surface:check}
\item \textbf{La prime de krach}~: la probabilité sous un niveau est la dérivée par rapport au strike du put, donc le skew
intervient par la dérivée de la volatilité.
\omcode{code/derivatives/07-implied-volatility-and-its-surface/python/dv_surface.py}{80}{95}{Probabilité risque-neutre d'une baisse, avec et sans skew.}\label{lst:dv:implied-volatility-and-its-surface:crash}
\item \textbf{Exécuter} \texttt{dv\_surface.check(index\_surface())} (aucune violation),
\texttt{check(index\_surface((1, -0.1, 0.03)))} (six), \texttt{crash\_probability()} et
\texttt{fig\_surface.py}.
\end{tutsteps}
\textbf{Pour aller plus loin.} Interpoler chaque smile linéairement en $k$ au lieu d'une spline cubique, et compter les violations
de butterfly que créent les cassures~; puis construire la surface sur une grille de strikes plutôt que sur une grille de
moneyness, et chercher les violations de calendrier au voisinage des dates de détachement.
\end{tutorial}

\section{Construction~: la surface de volatilité}\label{bld:dv:implied-volatility-and-its-surface:volsurface}

\begin{build}
\bfield{Rôle} L'objet unique de la petite entreprise pour la volatilité implicite~: les ajustements (chapitre~8),
les modèles (chapitres~9 à 13) et la bibliothèque de pricing (chapitre~28) le lisent, et rien en aval
ne doit voir d'arbitrage.
\bfield{Interface} \texttt{Surface(asof, expiries, ks, w, forwards)} avec
\texttt{implied\_vol(strike, expiry)}, \texttt{total\_variance(k, t)}, \texttt{forward(t)}~;
\texttt{from\_vols(asof, expiries, forwards, ks, vols)}~; \texttt{arbitrage\_report(surface)}~;
\texttt{density\_factor}~; \texttt{risk\_neutral\_density(strikes, calls, df)}.
\bfield{Règles} Interpoler la variance totale~: une spline cubique not-a-knot en $k$ au sein d'une échéance
(une spline naturelle infléchit à tort un smile convexe à ses extrémités), linéaire en $t$ à $k$ fixe
entre échéances, volatilité plate avant la première et après la dernière~; extrapoler linéairement en
$k$ et signaler les ailes comme un risque (chapitre~8).
\bfield{Tests d'acceptation} \texttt{code/firm/volsurface/tests/}~: la surface restitue ses volatilités d'entrée aux nœuds~;
une surface construite à partir d'une volatilité plate a un facteur de densité égal à un
en $k=0$ et aucune violation~; une cotation modifiée crée des violations de butterfly et une échéance longue
abaissée crée une violation de calendrier~; Breeden--Litzenberger retrouve la densité lognormale.
\bfield{Pour aller plus loin} Remplacer la grille de nœuds par les tranches SVI du chapitre~8, et ajouter un mode
sticky-delta qui recentre la surface sur un nouveau forward.
\end{build}

\begin{omsources}
\item E.~Derman, \emph{The Volatility Smile}, lecture notes, Columbia University (2008),
lecture 1~; et ``Regimes of volatility'', \emph{Risk} (avril 1999).
\item M.~Rubinstein, ``Implied binomial trees'', \emph{Journal of Finance} 49 (1994) 771--818.
\item D.~T.~Breeden et R.~H.~Litzenberger, ``Prices of state-contingent claims implicit in
option prices'', \emph{Journal of Business} 51 (1978) 621--651.
\item P.~Carr et D.~B.~Madan, ``A note on sufficient conditions for no arbitrage'',
\emph{Finance Research Letters} 2 (2005) 125--130.
\item G.~Bakshi, N.~Kapadia et D.~Madan, ``Stock return characteristics, skew laws, and the
differential pricing of individual equity options'', \emph{Review of Financial Studies} 16
(2003) 101--143.
\item Y.-F.~Chen et X.~Mu, ``Asymmetric volatility in commodity markets'', \emph{Journal of
Commodity Markets} 22 (2021).
\item J.~Gatheral, \emph{The Volatility Surface}, Wiley, 2006.
\end{omsources}

\section{Exercices}

\begin{exercise}[$\star$]\label{exo:dv:implied-volatility-and-its-surface:1}
Le forward à trois mois vaut 100.37. Donner la \omterm{def:dv:implied-volatility-and-its-surface:surface}{log-moneyness} du strike 90 et sa
variance totale pour une volatilité de 20.47.
\end{exercise}

\begin{exercise}[$\star$]\label{exo:dv:implied-volatility-and-its-surface:2}
Sur la surface du chapitre, les strikes 90 et 110 à trois mois sont à 20.47 et 14.19
de volatilité. Donner le skew 90--110 et expliquer pourquoi le même skew, exprimé en
termes de strike, serait plus faible à un an.
\end{exercise}

\begin{exercise}[$\star$]\label{exo:dv:implied-volatility-and-its-surface:3}
À même moneyness forward $k=0$, l'option à trois mois est cotée à 20~\% et
celle à six mois à 13~\%. Y a-t-il un \omterm{def:dv:implied-volatility-and-its-surface:static}{arbitrage calendaire}~? Quelle opération mettre en place~?
\end{exercise}

\begin{exercise}[$\star\star$]\label{exo:dv:implied-volatility-and-its-surface:4}
Calculer le prix du call digital à la monnaie à trois mois sur la surface du chapitre, avec et sans
la correction de skew, sachant une \omterm{def:dv:greeks-and-the-hedging-pnl:greeks}{vega} de 19.77 et une pente de volatilité $\partial\sigma/\partial K$
de $-0.0031$ par unité de strike.
\end{exercise}

\begin{exercise}[$\star\star$]\label{exo:dv:implied-volatility-and-its-surface:5}
Donner le delta du call à trois mois au strike 100 sous \omterm{def:dv:implied-volatility-and-its-surface:sticky}{sticky strike} et sous
\omterm{def:dv:implied-volatility-and-its-surface:sticky}{sticky delta}, et dire combien d'actions par tranche de 100 calls ce choix fait varier.
\end{exercise}

\begin{exercise}[$\star\star$]\label{exo:dv:implied-volatility-and-its-surface:6}
Pourquoi le skew d'un indice actions est-il plus pentu que ceux de ses composantes~?
\end{exercise}

\begin{exercise}[$\star\star\star$]\label{exo:dv:implied-volatility-and-its-surface:7}
\emph{Programmation.} Avec \texttt{index\_surface} et \texttt{check}, ajouter une erreur de cotation de $+3$
points à trois mois et $k=-0.1$~: combien de violations d'\omterm{def:dv:implied-volatility-and-its-surface:static}{arbitrage butterfly}, et de combien
$g$ devient négatif~? Trouver ensuite la plus petite erreur à la baisse sur la cotation à la monnaie à six mois qui
crée une violation d'\omterm{def:dv:implied-volatility-and-its-surface:static}{arbitrage calendaire}.
\end{exercise}

\begin{exercise}[$\star\star\star$]\label{exo:dv:implied-volatility-and-its-surface:8}
\emph{Trouver le défaut.} «~Le put 80 à trois mois est à 27 de volatilité contre 16 à la monnaie~:
le marché anticipe que l'indice sera deux fois plus volatil s'il descend à 80.~»
\end{exercise}

\section{Problème~: La prime de krach}

\begin{problem}[{Problème du week-end --- ce que dit le skew d'une baisse de 20~\%}]\label{pb:dv:implied-volatility-and-its-surface:1}
Un fonds de pension demande ce que disent les options sur indice d'une baisse de plus de 20~\% au cours des
trois prochains mois, et ce que coûte une protection contre elle. Utiliser la surface du chapitre~: spot 100,
taux 3~\%, rendement du dividende 1.5~\%, trois mois jusqu'à l'échéance.

\medskip\noindent\textbf{Partie I --- Les données.}
\begin{enumerate}
\item Donner le forward à trois mois et le facteur d'actualisation.
\item Donner la volatilité à la monnaie et la volatilité au strike 80.
\item Donner la \omterm{def:dv:implied-volatility-and-its-surface:surface}{log-moneyness} du strike 80.
\item Calculer le prix du put 80 sur la surface avec skew.
\item Calculer ce prix sur une surface plate à la volatilité à la monnaie.
\end{enumerate}

\medskip\noindent\textbf{Partie II --- Probabilités.}
\begin{enumerate}[resume]
\item Donner la probabilité risque-neutre de finir sous 80 sur la surface plate.
\item Donner la même probabilité sur la surface avec skew.
\item Pourquoi n'est-elle pas $\Phi(-d_2)$ à la volatilité propre du strike 80~?
\item Comment la probabilité se relie-t-elle à un digital put et à un put spread~?
\item Esquisser où le skew déplace la probabilité, à l'aide de la densité.
\end{enumerate}

\medskip\noindent\textbf{Partie III --- Interprétation.}
\begin{enumerate}[resume]
\item Le chiffre de 1.89~\% est-il une prévision de la fréquence de telles baisses~? Pourquoi pas~?
\item Qu'y a-t-il d'autre dans le prix du put 80~?
\item Comment une baisse de 5~\% de l'indice modifierait-elle le put 80 sous \omterm{def:dv:implied-volatility-and-its-surface:sticky}{sticky strike} et sous \omterm{def:dv:implied-volatility-and-its-surface:sticky}{sticky delta}~?
\item Pourquoi le skew est-il apparu après 1987~?
\item Pourquoi le skew pourrait-il être plus pentu à un mois qu'à un an~?
\end{enumerate}

\medskip\noindent\textbf{Partie IV --- Jugement.}
\begin{enumerate}[resume]
\item Que faudrait-il dire au fonds du coût de cette protection~?
\item Quel est le moyen le moins coûteux de se protéger contre une baisse de 20~\%, et à quoi renonce-t-il~?
\item Quelle surface un système de risque doit-il utiliser pour cette probabilité~: la surface plate ou la surface avec skew~?
\item Énoncer le \emph{résultat nommé}~: la probabilité risque-neutre d'une baisse sous 80 dans trois
mois sur la surface plate et sur la surface avec skew, ainsi que leur rapport.
\item En une phrase~: que valorise le skew~?
\end{enumerate}
\end{problem}

\section{Questions d'entretien}

\begin{interviewq}[$\star$ \iqroles{trader, chercheur}]\label{iq:dv:implied-volatility-and-its-surface:1}
Pourquoi les options sur indices actions ont-elles un skew, et pourquoi n'est-il apparu qu'après 1987~?
\end{interviewq}

\begin{interviewq}[$\star$ \iqroles{chercheur}]\label{iq:dv:implied-volatility-and-its-surface:2}
Comment obtenez-vous la \omterm{def:dv:implied-volatility-and-its-surface:density}{densité risque-neutre} à partir des prix d'options~?
\end{interviewq}

\begin{interviewq}[$\star\star$ \iqroles{trader}]\label{iq:dv:implied-volatility-and-its-surface:3}
L'indice baisse de 5~\%. Sous \omterm{def:dv:implied-volatility-and-its-surface:sticky}{sticky strike} et sous
\omterm{def:dv:implied-volatility-and-its-surface:sticky}{sticky delta}, que deviennent la volatilité à la monnaie et votre delta~?
\end{interviewq}

\begin{interviewq}[$\star\star$ \iqroles{chercheur, développeur}]\label{iq:dv:implied-volatility-and-its-surface:4}
Donnez les conditions statiques d'absence d'arbitrage sur une \omterm{def:dv:implied-volatility-and-its-surface:surface}{surface de volatilité}
implicite, et comment vous les vérifiez.
\end{interviewq}

\begin{interviewq}[$\star\star$ \iqroles{risque}]\label{iq:dv:implied-volatility-and-its-surface:5}
Un rapport de risque calcule la probabilité d'une baisse de 20~\% à partir de la volatilité à la monnaie. Qu'est-ce
qui ne va pas, et de combien peut-il se tromper~?
\end{interviewq}

\begin{interviewq}[$\star\star\star$ \iqroles{développeur, chercheur}]\label{iq:dv:implied-volatility-and-its-surface:6}
Comment interpoleriez-vous une surface entre les strikes et les échéances cotés, sans créer
d'arbitrage~?
\end{interviewq}
